\documentclass[10pt,a4paper]{article}
\usepackage[T1]{fontenc}
\usepackage[utf8]{inputenc}
\usepackage{lmodern,textcomp}
\usepackage[a4paper,hmargin=2.1cm,vmargin=2cm]{geometry}
\usepackage{amsmath,amssymb}
\usepackage{microtype}
\usepackage[table]{xcolor}
\usepackage{booktabs,tabularx,array}
\usepackage{graphicx}
\usepackage{float}
\usepackage[font=small,labelfont=bf]{caption}
\usepackage{enumitem}
\usepackage[hidelinks]{hyperref}
\setlist{itemsep=1pt,topsep=3pt}
\renewcommand{\topfraction}{0.9}
\renewcommand{\bottomfraction}{0.8}
\renewcommand{\textfraction}{0.07}
\renewcommand{\floatpagefraction}{0.85}
\setcounter{totalnumber}{4}
\renewcommand{\arraystretch}{1.3}
\renewcommand{\tabularxcolumn}[1]{m{#1}}
\setlength{\aboverulesep}{0pt}
\setlength{\belowrulesep}{0pt}
\definecolor{zebra}{gray}{0.93}
\newcolumntype{L}{>{\raggedright\arraybackslash}X}
\title{LOGOS: Benchmarks for DDA Peptide Search}
\author{Filip Rumenovski}
\date{1 November 2026}
\begin{document}
\maketitle
\begin{abstract}
LOGOS is a proteomics search engine written in Rust with a CUDA backend. It reads Thermo RAW files directly and searches on one CPU node or one GPU. Against Sage 0.14.7 on the same node, chemistry and 22 cores, it runs closed DDA 12.2$\times$ and open search 6.6$\times$ faster on the CPU, and open search 16.6$\times$ faster on one GPU, while reporting more peptide sequences at an entrapment-measured error. At scale it searches the 79.7 GB HEK293 deposit PXD001468 in an open search in 12.13 s on one GPU, and one gut metaproteome against 111.5 million sequences in 20.3 s on an L40S (recorded on an earlier build). CPU and GPU builds produce byte-identical identifications. The record behind this report covers 216 RAW files from 11 public datasets and 172 recorded comparisons against external tools, released with it as data.
Measurements were taken through 31 October 2026 on the Warrior cluster at Wayne State University; this report covers data-dependent acquisition only.
\end{abstract}

\section{Across datasets}
Figure~\ref{fig:force} sets LOGOS against the strongest tool or published result available on each benchmark: Sage where Sage was run on the same files and node, the authors' deposited search where that is the reference, and CHIMERYS, a chimeric-spectrum engine, where its deposits exist. Counts are distinct sequences or PSMs at 1\% FDR under one counting rule; a filled marker means LOGOS's count is the entrapment-validated one. Table~\ref{tab:board} gives the numbers behind every bar.

\begin{figure}[H]\centering
\includegraphics[width=\linewidth]{fig-force.pdf}\par
\caption{(a) LOGOS identifications relative to the strongest comparator on 13 benchmarks across 6 public datasets, 4 instrument families and 9 search tasks. Dark bars: ahead by more than 1\%; light: within 1\%; hatched: behind. Filled marker: entrapment-validated LOGOS count; open: target-decoy only. (b) Wall-time speed-up over Sage; the whisker spans the three nodes measured for open search (Section~\ref{sec:sage}). Rows marked GPU compare one GPU with Sage on CPU.}\label{fig:force}
\end{figure}

\begin{table}[H]\centering
\caption{Scoreboard: LOGOS against the strongest comparator on each benchmark. Status gives the date of the measurement; rows from 31 October were rerun on the current build.}\label{tab:board}
\scriptsize\setlength{\tabcolsep}{2.6pt}
\begin{tabularx}{\linewidth}{@{}>{\raggedright\arraybackslash}m{1.15in} >{\raggedright\arraybackslash}m{0.98in} >{\raggedright\arraybackslash}m{0.86in} rr r L >{\raggedright\arraybackslash}m{0.62in}@{}}
\toprule
Benchmark & Data & Comparator & LOGOS & Comparator & $\Delta$ & Error control & Status\\
\midrule
\rowcolor{zebra}LFQ human, open \ensuremath{-}100 to +500 Da & PXD028735\newline Q Exactive HF-X & Sage 0.14.7 & 35,595 & 30,343 & \textbf{+17.3\%} & entrapment 1.04\% vs 0.90\% & measured 2026-10-31\\
LFQ human, closed & PXD028735\newline Q Exactive HF-X & Sage 0.14.7 & 31,718 & 31,489 & {+0.7\%} & entrapment 1.24\% each & measured 2026-10-30\\
\rowcolor{zebra}Astral TMTpro, one file & PXD058160\newline Orbitrap Astral & Sage 0.14.7 & 7,568 & 7,212 & \textbf{+4.9\%} & entrapment 0.76\% vs 1.02\% & measured 2026-10-31\\
Astral TMTpro, 48 files & PXD058160\newline Orbitrap Astral & Authors' Comet + LDA & 1,700,508 & 1,333,804 & \textbf{+27.5\%} & entrapment 0.997\%; ahead in 48 of 48 files & recorded 2026-10-22\\
\rowcolor{zebra}HLA class I, 9 alleles & MSV000084172\newline Orbitrap & Deposited Spectrum Mill & 20,625 & 17,248 & \textbf{+19.6\%} & target-decoy; reference without matched FDR & recorded 2026-09-19\\
Astral HEK, 30-min gradient & PXD046453\newline Orbitrap Astral & CHIMERYS (recounted) & 91,865 & 91,362 & {+0.6\%} & LOGOS entrapment 1.02\% & recorded 2026-09-29\\
\rowcolor{zebra}Astral HEK, 14-min gradient & PXD046453\newline Orbitrap Astral & CHIMERYS (recounted) & 55,383 & 55,309 & {+0.1\%} & LOGOS entrapment 1.03\% & recorded 2026-09-29\\
Mouse, 1 Th isolation & PXD053241\newline Q Exactive HF-X & CHIMERYS (recounted) & 68,896 & 68,624 & {+0.4\%} & LOGOS entrapment 1.04\% & recorded 2026-09-29\\
\rowcolor{zebra}Mouse, 20 Th isolation & PXD053241\newline Q Exactive HF-X & CHIMERYS (recounted) & 39,971 & 40,007 & {-0.1\%} & LOGOS entrapment 1.05\% & recorded 2026-09-24\\
Ion trap CID, 0.4 Th & PXD053241\newline Orbitrap Eclipse & CHIMERYS (recounted) & 21,364 & 20,941 & \textbf{+2.0\%} & LOGOS entrapment 1.00\%; CHIMERYS's database & recorded 2026-09-29\\
\rowcolor{zebra}LFQ human, chimeric & PXD053241\newline Q Exactive HF-X & CHIMERYS (recounted) & 59,606 & 61,223 & {-2.6\%} & LOGOS entrapment 1.05\% & recorded 2026-09-23\\
Ion trap HCD, 0.4 Th & PXD053241\newline Orbitrap Eclipse & CHIMERYS (recounted) & 46,340 & 50,624 & {-8.5\%} & LOGOS entrapment 0.98\% & recorded 2026-09-29\\
\rowcolor{zebra}PTM: acetyl, phospho, ubiquitin & PXD053241, PXD013615\newline Q Exactive HF-X & CHIMERYS (recounted) & 41,985 & 44,560 & {-5.8\%} & target-decoy & recorded 2026-09-23\\
\bottomrule
\end{tabularx}
\end{table}

\textbf{Reading the scoreboard.} Against Sage LOGOS reports more sequences on every dataset tested: 0.7\% more in closed search at equal measured error, 4.9\% more on Astral TMTpro and 17.3\% more in open search, while running 12.2$\times$ to 16.6$\times$ faster. Against the authors' own searches it reports 27.5\% more PSMs on a 48-file TMTpro study, ahead in all 48 files, and 19.6\% more peptides on HLA immunopeptidomics. Against CHIMERYS it ties within 1\% on Astral gradients of 14 and 30 minutes and on wide-window data at 1 and 20 Th, edges ahead on ion-trap CID, and trails on ion-trap HCD, PTM sets and chimeric LFQ.

\section{Speed and scale}\label{sec:scale}
Figure~\ref{fig:power} collects every speed comparison in the record in which LOGOS and another tool did the same work on the same input. The search rows are the Sage comparisons of Section~\ref{sec:sage}. Rescoring is 11--34$\times$ faster than C++ Percolator on the same input, and the MSBooster features are 75$\times$ faster with byte-identical output. LOGOS searches a RAW file to identifications 94$\times$ faster than ThermoRawFileParser converts that file to the mzML other engines read.

\begin{figure}[H]\centering
\includegraphics[width=\linewidth]{fig-power.pdf}\par
\caption{Speed-up of LOGOS over each comparator, as the comparator's time divided by LOGOS's, on the same input and work. Search: wall time against Sage 0.14.7 on CPU on the same node, LOGOS on CPU (dark) or on one GPU (light). Rescoring and features: CPU, same PIN input; MSBooster's 58.29\,s includes about 26\,s of Java start-up. Whole pipeline: LOGOS's full search of one Astral file from RAW against ThermoRawFileParser converting that file alone; LOGOS's median search call per MS2 (search only, $\pm$50\,ppm precursor window, CPU, 4,096 MS2 of PXD017823) against Orbiter's published 17\,ms median, on different hardware.}\label{fig:power}
\end{figure}

Where no other engine has been run on the same input, Table~\ref{tab:scale} reports absolute throughput. LOGOS searches the 24-file HEK293 deposit of Chick et al.\ (PXD001468, 79.70\,GB) in an open search in 12.13\,s on one GPU, reporting 132,592 sequences at an entrapment-measured 1.03\%. It searches one gut-metaproteome run against the 111.5 million sequences of ComPIL~2.0, scoring 178.9 billion candidates, in 20.26\,s on an L40S as recorded on 18 September; the current build takes 65.05\,s on an RTX PRO 4500 with identical identifications.

\begin{table}[H]\centering
\caption{Throughput at scale and latency in real time. Inputs staged in RAM; GPU counts are physical devices. Rows were recorded on earlier builds before 31 October 2026 and not rerun on the current build.}\label{tab:scale}
\scriptsize\setlength{\tabcolsep}{3pt}
\begin{tabularx}{\linewidth}{@{}>{\raggedright\arraybackslash}m{1.75in} >{\raggedright\arraybackslash}m{1.2in} L >{\raggedright\arraybackslash}m{0.95in} >{\raggedright\arraybackslash}m{0.85in}@{}}
\toprule
Workload & Result & Identifications & Hardware & Note\\
\midrule
\rowcolor{zebra}HEK293 open search, PXD001468: 24 RAW, 79.70 GB, 1,121,149 MS2 & 12.13 s; 394 GB/\allowbreak{}min; 92,428 MS2/\allowbreak{}s & 132,592 sequences, 688,428 PSMs at 1\%; entrapment FDP 1.03\% & 1x H200 NVL, 32 EPYC 9554 cores & median of 3; RAWs in RAM\\
Same files, closed search & 7.89 s; 606 GB/\allowbreak{}min & -- & 1x H200 NVL & RAWs in RAM\\
\rowcolor{zebra}Same open search on a workstation GPU & 20.30 s; 236 GB/\allowbreak{}min & -- & 1x RTX PRO 4000, 32 EPYC 9454 cores & RAWs in RAM\\
HEK293 cascade: every MS2 through nine layers & 939,455 of 1,121,149 MS2 explained (83.8\%) & entrapment FDP 0.87\% overall (Ne 4,025) & 1x L40S & two layers above 1\% alone; without them 81.96\%\\
\rowcolor{zebra}ComPIL 2.0 metaproteome: one gut RAW against 111,478,540 sequences (84 GB FASTA, 64.06 G residues) & 20.26 s; 178.9 billion candidates scored & 38,649 PSMs, 10,326 peptides, 8,077 proteins at 1\% (target-decoy) & 1x L40S, 32 EPYC 9554 cores & recorded 2026-09-18, median of 2; performance settings; no entrapment arm\\
ComPIL 2.0, same search on the current build & 65.05 s (repeats 65.05, 65.11, 64.96); same 178.9 billion candidates & identical: 38,649 PSMs, 10,326 peptides, 8,077 proteins & 1x NVIDIA RTX PRO 4500 Blackwell, 32 EPYC 9454 cores & run 2026-10-30-compil2; GPU-bound: the scan is 62.9 s of the run\\
\rowcolor{zebra}ComPIL 2.0, published-protocol settings & 80.2\% of the 7,565 published H2 peptides recovered & after I/\allowbreak{}L normalisation & Apple GPU (Metal) & separate run from the speed run\\
Astral TMTpro, PXD058160: 15 RAW, 39.5 GB & 20.91 s; 113 GB/\allowbreak{}min & 392,305 PSMs, 97,666 sequences at 1\% & 1x RTX PRO 4000 & median of 3; one process\\
\bottomrule
\end{tabularx}
\end{table}

The cascade asks what share of every MS2 spectrum in PXD001468 can be explained at controlled error. Each layer searches only the spectra that every earlier layer left unexplained. One learner scores the rows of all layers, and target-decoy competition runs per spectrum across layers. The error is measured on a separate entrapment arm that runs the whole layered pipeline. Figure~\ref{fig:cascade} shows the result: 939,455 of 1,121,149 spectra (83.8\%) explained at an overall entrapment FDP of 0.87\%.

\begin{figure}[H]\centering
\includegraphics[width=\linewidth]{fig-cascade.pdf}\par
\caption{The HEK293 cascade on all 24 files of PXD001468. Each bar is the share of all MS2 that a layer explains, placed after the layers before it; the label gives spectra explained and that layer's own entrapment FDP. Hatched layers run above 1\% on their own (1.04--1.39\%); without the two highest, Unimod offsets and missed cleavage, the cascade explains 918,865 spectra (82.0\%). Of the spectra left unexplained, about half are chimeric and a third carry too little information for any engine; well-formed spectra that stay unexplained are 3.3\% of the total (6-file subset).}\label{fig:cascade}
\end{figure}

\section{Real-time acquisition decisions}\label{sec:rt}
In real-time search the instrument asks, for every MS2 it records, whether to spend an MS3 scan on that precursor. LOGOS answers with one pipeline per MS2: candidate search, XCorr, a lookup of predicted fragment spectra, a causal retention-time check, a frozen score model and a trigger policy. The measurements below replay recorded acquisitions of PXD043660 (Thermo real-time search on Orbitrap Ascend and Eclipse) and PXD058160 (Orbitrap Astral) one MS2 at a time, at batch 1, with RAW decoding excluded; every percentile is recomputed from the replay's raw per-MS2 samples by \texttt{eval/latency.py}.

On the TKO file a decision takes a median 0.884\,ms and 1.205\,ms at the 99th percentile on one RTX PRO 4500, over three interleaved runs of 27,674 decisions, every one of which equals the offline reference. The previous binary measured 0.925\,ms in the same job. One of the 83,022 timed decisions exceeded the 80\,ms budget Thermo's implementation sets, a single 11.3\,s host stall; on the HYpro file 19 of 96,600 did. On the Astral file, whose recorded MS2 arrive a median 24.8\,ms apart, the median decision takes 1.98\,ms and none of 135,711 exceeds 80\,ms, from a single run.

\begin{figure}[H]\centering
\includegraphics[width=\linewidth]{fig-realtime.pdf}\par
\caption{(a) The median time of each stage of a TKO decision; the stage medians sum to 0.87\,ms, slightly less than the median whole decision, as medians do. (b) Each replay's median (filled) and 99th-percentile (open) decision time against its own budget (bar): the median time to the next MS2 for the Orbitrap Ascend files, the median recorded MS2 gap for the Astral file. The dashed line is the 80\,ms cap of Thermo's real-time search.}\label{fig:rt}
\end{figure}

\begin{table}[H]\centering
\caption{Decision latency, recomputed from raw per-MS2 samples (runs \texttt{2026-10-22-rts-latency}, \texttt{2026-10-22-astral-latency}; job 40703237, RTX PRO 4500, EPYC 9454). Times in ms; over-80\,ms counts are over all runs.}\label{tab:rtlat}
\scriptsize\setlength{\tabcolsep}{3pt}
\begin{tabularx}{\linewidth}{@{}L l r r r r >{\raggedright\arraybackslash}m{1.05in} r@{}}
\toprule
Replay & Binary & Runs & p50 & p99 & Over 80\,ms & Matching the offline reference & Budget\\
\midrule
\rowcolor{zebra}TKO, Orbitrap Ascend (27,674 MS2) & s7 & 3 & 0.884 & 1.205 & 1 of 83,022 & 27,674 of 27,674 & 81.0\\
HYpro, Orbitrap Ascend (32,200 MS2) & s7 & 3 & 1.237 & 1.868 & 19 of 96,600 & 32,200 of 32,200 & 73.8\\
\rowcolor{zebra}Astral TMTpro (135,711 MS2) & astral-opt & 1 & 1.983 & 3.015 & 0 of 135,711 & 135,704 exact of 135,711 & 24.8\\
\bottomrule
\end{tabularx}
\end{table}

Speed matters only if the decisions are good. Replaying all 24 PXD043660 acquisitions with the authors' own database, and trigger rules chosen on files other than the one being replayed, LOGOS triggers on more of the authors' passing scans with fewer MS3 scans, and the scans it triggers carry an entrapment-measured FDP of 0.56\% where Thermo's recorded triggers carry 2.95\%, lower in all 24 files (sign test $p = 1.2\times10^{-7}$). It does not win everywhere: raw protein counts tie at 5,442, which the rule we set ourselves counts as a loss, both HYpro regimes trail by up to 10 proteins, and LOGOS has more proteins in only 6 of the 24 files. This replay and the latency runs use different recipes (the authors' database and v11 rules here; an earlier database and v8 rules for the timing), so neither number is claimed for the other.

\begin{table}[H]\centering
\caption{Replay of 24 Thermo real-time-search acquisitions (PXD043660), LOGOS / Thermo recorded; bold where LOGOS is higher. Entrapment FDP in percent, of the triggered scans' identifications, both masks classified by the same scorer. Recorded 21 October 2026 (run \texttt{2026-10-21-rts-replay}).}\label{tab:rtreplay}
\scriptsize\setlength{\tabcolsep}{3pt}
\begin{tabularx}{\linewidth}{@{}L r r r r r@{}}
\toprule
Regime & Author-pass scans & Modified peptides & Raw proteins & MS3 scans & Entrapment FDP\\
\midrule
\rowcolor{zebra}HYpro, Ascend (2 files) & \textbf{15,954} / 15,950 & \textbf{9,542} / 9,390 & 2,322 / 2,332 & 20,036 / 20,041 & \textbf{1.79} / 6.97\\
HYpro, Eclipse (2 files) & \textbf{17,408} / 17,194 & \textbf{10,581} / 10,473 & 2,673 / 2,675 & 20,333 / 20,564 & \textbf{1.48} / 6.20\\
\rowcolor{zebra}TKO11, Ascend (10 files) & \textbf{67,328} / 67,176 & \textbf{16,259} / 15,973 & \textbf{2,528} / 2,518 & 73,561 / 73,645 & \textbf{0.27} / 1.75\\
TKO11, Eclipse (10 files) & \textbf{67,024} / 66,693 & \textbf{16,545} / 16,229 & \textbf{2,665} / 2,655 & 72,999 / 73,202 & \textbf{0.24} / 2.21\\
\rowcolor{zebra}Pooled, 24 files & \textbf{167,714} / 167,013 & \textbf{30,766} / 30,182 & 5,442 / 5,442 & 186,929 / 187,452 & \textbf{0.56} / 2.95\\
\bottomrule
\end{tabularx}
\end{table}

On the 12 Orbiter acquisitions of PXD017823 the trigger rules were chosen retrospectively and the database differs from the authors', so these results are weaker evidence. One rule triggers 38,169 of the authors' passing scans against Orbiter's 30,815, with 40,340 MS3 scans against 41,647, at an estimated 2.7\% false-trigger rate against 16.4\%, but reaches fewer raw protein accessions (8,857 against 9,245). Another reaches 8,710 proteins at 1\% picked-protein FDR against 8,315, 6,563 of them with two or more peptides against 5,821, using 39,536 MS3 scans. A separate measurement times the search call alone with a broad closed precursor window of $\pm$500\,Da: 32.7\,ms per MS2 at the median on an NVIDIA A30, against 2.74\,s on CPU.

\section{Head to head with Sage on the current build}\label{sec:sage}
Both engines ran in the same Slurm job on the same node, interleaved with the order rotated between repeats; times are medians of three or more wall-clock measurements from input file to identifications, container start included. Sage's speed varies by a factor of 2.8 between these nodes, so every ratio is taken within one node.

\begin{table}[H]\centering
\caption{Wall time per search (seconds, median) and speed relative to Sage on the same node.}\label{tab:speed}
\footnotesize\setlength{\tabcolsep}{4pt}
\begin{tabularx}{\linewidth}{@{}>{\raggedright\arraybackslash}m{1.35in}>{\raggedright\arraybackslash}m{1.18in} rrr rr lL@{}}
\toprule
& & \multicolumn{3}{c}{Wall time (s)} & \multicolumn{2}{c}{Speed-up} & & \\
\cmidrule(lr){3-5}\cmidrule(lr){6-7}
Benchmark & Node & CPU & GPU & Sage & CPU & GPU & Binary & Slurm job\\
\midrule
\rowcolor{zebra}Closed DDA, LFQ human & ymh4 \textperiodcentered{} EPYC 9454 & 1.41 & 1.53 & 17.25 & 12.2 & 11.3 & 47068225 & 40931267\\
Open \ensuremath{-}100 to +500 Da, LFQ human & ymh4 \textperiodcentered{} EPYC 9454 & 6.96 & 2.79 & 46.23 & 6.6 & 16.6 & 47068225 & 40928343\\
\rowcolor{zebra}Open \ensuremath{-}100 to +500 Da, LFQ human & bbc1 \textperiodcentered{} EPYC 9374F & 7.08 & -- & 37.79 & 5.3 & -- & 47068225 & 40928337\\
Open \ensuremath{-}100 to +500 Da, LFQ human & arw3 \textperiodcentered{} Gold 6326 & 11.36 & -- & 106.4 & 9.4 & -- & 6671172e & 40928265\\
\rowcolor{zebra}Astral TMTpro, one file & ymh4 \textperiodcentered{} EPYC 9454 & 3.68 & 2.58 & 40.06 & 10.9 & 15.5 & 47068225 & 40931267\\
\bottomrule
\end{tabularx}
\end{table}

Identification counts use one rule for both engines: Sage's results are converted to the same input format and counted by the same program, the open counter \texttt{eval/count.py}, which gives exactly the engine's integers on every counted run (run \texttt{2026-11-01-eval-parity}). The entrapment arm adds a second database, equal in size and absent from the sample, to each search; every accepted hit to it is a known false discovery, and $\mathrm{FDP} = N_e(1 + 1/r)/(N_t + N_e)$, with the ratio $r$ of entrapment to target sequence space computed from each search's own digestion (1.0606 for the LFQ human database, 0.9998 for the Astral mouse database).

\begin{table}[H]\centering
\caption{Identifications at 1\% FDR on the target database alone and in the entrapment arm, at sequence and PSM level. FDP in percent.}\label{tab:ids}
\footnotesize\setlength{\tabcolsep}{2.6pt}
\begin{tabularx}{\linewidth}{@{}L rr c rrr c rrr@{}}
\toprule
& \multicolumn{2}{c}{Target only} && \multicolumn{3}{c}{Entrapment arm, sequences} && \multicolumn{3}{c}{Entrapment arm, PSMs}\\
\cmidrule{2-3}\cmidrule{5-7}\cmidrule{9-11}
Engine & Sequences & PSMs && $N_t$ & $N_e$ & FDP && $N_t$ & $N_e$ & FDP\\
\midrule
\multicolumn{10}{@{}l}{\textit{Closed DDA, LFQ human}}\\
\rowcolor{zebra}\quad LOGOS & 33,102 & 59,769 && 31,718 & 204 & 1.24 && 57,154 & 349 & 1.18\\
\quad Sage 0.14.7 & 32,453 & 58,668 && 31,489 & 202 & 1.24 && 56,594 & 337 & 1.15\\
\addlinespace[2pt]
\multicolumn{10}{@{}l}{\textit{Open \ensuremath{-}100 to +500 Da, LFQ human}}\\
\rowcolor{zebra}\quad LOGOS & 34,130 & 70,886 && 32,086 & 166 & 1.00 && 65,824 & 353 & 1.04\\
\quad LOGOS + LDA rescoring & -- & -- && 35,650 & 203 & 1.10 && 73,449 & 390 & 1.03\\
\rowcolor{zebra}\quad LOGOS + calibration + LDA & 37,178 & 77,433 && 35,595 & 192 & 1.04 && 73,438 & 379 & 1.00\\
\quad Sage 0.14.7 & 31,744 & 73,148 && 30,343 & 141 & 0.90 && 68,848 & 359 & 1.01\\
\addlinespace[2pt]
\multicolumn{10}{@{}l}{\textit{Astral TMTpro, one file}}\\
\rowcolor{zebra}\quad LOGOS & 7,769 & 24,167 && 7,568 & 29 & 0.76 && 23,364 & 124 & 1.06\\
\quad Sage 0.14.7 & 7,470 & 22,409 && 7,212 & 37 & 1.02 && 21,649 & 94 & 0.86\\
\bottomrule
\end{tabularx}
\end{table}

LDA rescoring alone raises the open-search error to 1.10\%; with mass calibration it holds 1.04\%, the configuration reported. On Astral TMTpro LOGOS reports more sequences at lower error, while at the PSM level Sage is the more conservative.

\begin{figure}[H]\centering
\includegraphics[width=\linewidth]{fig-throughput.pdf}\par
\caption{(a) Wall time per file when several Astral TMTpro files are searched in one run; each file's output is byte-identical to that file searched alone (jobs 40912674, 40912675, 40911945). (b) ThermoRawFileParser 2.0.0-dev converting one 2.6 GB Astral file to the mzML that Sage reads, alone and six at once (job 40913997), against LOGOS's whole search of that file from RAW.}\label{fig:thru}
\end{figure}

\section{Method and correctness}
\begin{itemize}
\item \textbf{Hardware.} Warrior compute nodes at Wayne State University, shared with other tenants: ymh4 (AMD EPYC 9454, NVIDIA RTX PRO 4500 Blackwell), bbc1 (AMD EPYC 9374F), arw3 (Intel Xeon Gold 6326); 22 cores per job, inputs staged in RAM.
\item \textbf{Chemistry.} LFQ human: trypsin, one missed cleavage, 7--30 residues, carbamidomethyl C, up to two oxidised M, 10 ppm precursor (open: $-100$ to $+500$ Da), 20 ppm fragments, charges 2--4, deisotoping. Astral TMTpro: TMTpro on K and the N-terminus, two missed cleavages, 7--50 residues, 50 ppm precursor, 10 ppm fragments, charges 2--6. Sage 0.14.7 release defaults with this chemistry, including its LDA rescoring and retention-time prediction.
\item \textbf{Backend parity and determinism.} CPU and GPU builds produce byte-identical identification files (open search: the same hash, 583ae51be609\dots, on every run); the same binary, input and settings give the same output on every run, because reductions are order-independent.
\item \textbf{The CPU path is the reference.} GPU kernels are tested against it, and new CPU code against brute-force implementations, mutation-checked so a broken variant fails its test.
\item \textbf{Evidence.} Every number measured on the current build has a run record (\texttt{runs/}) naming the node, the binary and its SHA-256, the SHA-256 of every input, every repeat's wall time and the hash of every output count; the speed and identification tables of this report are generated from those records and a check fails the build if they differ. Inputs are pinned to their public deposits by SHA-256, the entrapment database is regenerated bit-for-bit by its generator, and the counting rule is open code with tests.
\item \textbf{What LOGOS takes from Sage.} The CPU open-search preselection uses Sage's fragment-index layout (mass-sorted buckets of 8,192 peptides, fragments sorted within each); Sage's LFQ code was the design reference for LOGOS's MS1 quantification. On Astral TMTpro Sage spends 23.3 of its 39.4 s reading mzML at about four busy cores, so LOGOS's lead there comes mostly from reading RAW and preparing the database rather than the search core alone.
\end{itemize}

\section{Limitations}
LOGOS trails CHIMERYS on ion-trap HCD ($-8.5$\%), on three PTM sets ($-5.8$\% combined) and on chimeric LFQ ($-2.6$\%), and quantifies half as many precursors as CHIMERYS on the LFQbench mixtures at equal accuracy (17,762 against 35,053; median error 0.175 against 0.177). Rows dated before 31 October come from earlier builds and are reported as recorded; several compare against a recount of the reference's deposited results rather than a fresh run. Speed ratios depend on the node, and every node was shared with other jobs. The complete record behind this report, 172 comparisons with their caveats and notebook sources, is in \texttt{reports/data/catalog.json}.

\section{Scope of the record}
Table~\ref{tab:cover} lists the public datasets behind the benchmark record and the tools each was compared against.

\begin{table}[H]\centering
\caption{Public datasets in the benchmark record. Rows counts the recorded measurements that cite the accession.}\label{tab:cover}
\scriptsize\setlength{\tabcolsep}{3pt}
\begin{tabularx}{\linewidth}{@{}>{\raggedright\arraybackslash}m{1.05in} >{\raggedright\arraybackslash}m{1.2in} >{\raggedright\arraybackslash}m{0.75in} r L >{\raggedright\arraybackslash}m{1.2in} r@{}}
\toprule
Accession & Dataset & Instrument & RAW & Benchmarked & Compared against & Rows\\
\midrule
\rowcolor{zebra}PXD001468 & HEK293 proteome, Chick et al. 2015 & Q Exactive & 24 & Open search at scale, closed search, explaining every spectrum in layers & Published counts & 154\\
PXD058160 & Astral TMTpro, four mouse tissues & Orbitrap Astral & 48 & TMT identification depth, speed vs Sage, multi-file throughput & Authors' Comet + LDA deposit; Sage 0.14.7 & 31\\
\rowcolor{zebra}PXD043660 + PXD023203 + PXD017803 & Thermo real-time search runs, TMT SPS-MS3 & Orbitrap Eclipse and Ascend & 32 & Real-time acquisition decisions replayed, error rate of triggered MS3 & Thermo's recorded real-time search triggers & 73\\
PXD017823 & Real-time search TMT study & Orbitrap & 24 & Identification depth and protein-level error & Authors' deposited passing scans & 42\\
\rowcolor{zebra}PXD058918 & Astral TMTpro18 fractions & Orbitrap Astral & 24 & TMTpro18 identification depth & Authors' deposited search & 9\\
PXD046453 & Astral DDA, HEK 1 \ensuremath{\mu}g, 14/15/30-min gradients & Orbitrap Astral & 15 & Chimeric DDA+ depth on Astral 14- and 30-minute gradients & CHIMERYS (recounted) & 107\\
\rowcolor{zebra}PXD053241 (+ PXD023218, PXD013615) & CHIMERYS benchmark deposits: LFQ human, wide-window mouse, ion trap, PTM & Q Exactive HF-X, Orbitrap Eclipse & 21 & Wide-window chimeric DDA (8 widths), ion-trap HCD and CID (9), PTM sets (3), LFQ human & CHIMERYS, Sequest HT, Comet, MS-GF+, MetaMorpheus deposits & 72\\
PXD028735 & LFQbench three-species mixtures and LFQ human & Q Exactive HF-X & 7 & Closed and open search speed and depth vs Sage; quantification accuracy and depth & Sage 0.14.7, CHIMERYS, ProteoBench submissions & 31\\
\rowcolor{zebra}MSV000084172 & HLA immunopeptidomics, 9 alleles & Orbitrap & 18 & Non-specific peptide identification & Deposited Spectrum Mill results & 4\\
PXD022433 / MSV000082943 & ComPIL 2.0 gut metaproteome & Orbitrap & 1 & Search against 111.5 M protein sequences & Published ProLuCID / DTASelect peptides & 15\\
\rowcolor{zebra}PXD012986, PXD000612 & Mouse tissue; HeLa phospho & Orbitrap & 2 & Rescoring models & Percolator & 10\\
\midrule
\rowcolor{zebra}Total & & & 216 & & & \\
\bottomrule
\end{tabularx}
\end{table}

\end{document}
